\exnumber{03}
\extitle{UI IV - Ex 4 - Animation}
\exfiles{ex03/fichier Figma, nomme Name\_FirstName\_DeliverableName\_Date\_VersionNumber}
\exauthorize{Ceux nommés dans l'exercice}
\makeheaderfiles

\textbf{Instructions :}

Animez vos 5 écrans sur Figma. Pensez au multistate !!

\textbf{Livrable :}

\begin{itemize}
\item Un prototype haute fidélité animé sur Figma
\item Déposez le fichier Figma sur Github, nommé ainsi : « Name\_FirstName\_DeliverableName\_Date\_VersionNumber ». Exemple : « Doe\_John\_MidFiPrototype\_05242024\_V1 ».
\item Ajoutez un titre dans le fichier.
\item Soignez l'aspect global du fichier. Il doit être clair, simple et lisible. Vous pouvez regarder canva.com, sans surcharger le dessin de détails.
\end{itemize}

\textbf{Ressources :}

\begin{itemize}
\item Gérer efficacement vos designs - Comment Mockplus State Page aide à voir et gérer plusieurs états \url{https://help.mockplus.com/p/372}
\end{itemize}

\textbf{Citation :}

« Parlez humain : utilisez les mots que vos utilisateurs utilisent. » Dan Saffer | Microinteractions

